\newcommand{\makeUC}[9]{
    \ifblank{#1}{}
    {
        \begin{figure}[H]
            \centering
            \includegraphics[width=0.9\textwidth]{usecase/#1.png}
            \caption{Diagramma #1}
        \end{figure}
    }
    \begin{table}[H]
        \centering
        \small
        \renewcommand{\arraystretch}{1.4}
        \begin{tabular}{|p{3cm}|p{10cm}|}
            \hline
            \textbf{Attori} & #2 \\
            \hline
            \textbf{Pre-condizioni} & #3 \\
            \hline
            \textbf{Post-condizioni} & #4 \\
            \hline
            \textbf{Scenario principale} & #5 \\
            \hline
            \ifblank{#6}{}{
            \textbf{Inclusioni} & #6 \\
            \hline
            }
            \ifblank{#7}{}{
            \textbf{Scenari alternativi} & #7 \\
            \hline
            }
            \ifblank{#8}{}{
            \textbf{Estensioni} & #8 \\
            \hline
            }
            \ifblank{#9}{}{
            \textbf{Generalizzazioni} & #9 \\
            \hline
            }
        \end{tabular}
    \end{table}
}

% contatore RF
\newcounter{rcount}
\setcounter{rcount}{0}

\newcommand{\makeReq}[3]{%
    \stepcounter{rcount}%
    RF-\thercount & 
    \ifcase#1\or Obbligatorio\or Opzionale\or Desiderabile\fi & 
    #2 & 
    UC#3 §\ref{#3} \\
    \hline
}

% contatore RQ
\newcounter{rqcount}
\setcounter{rqcount}{0}

\newcommand{\RQ}[3]{%
    \stepcounter{rqcount}%
    RQ-\therqcount & 
    \ifcase#1\or Obbligatorio\or Opzionale\or Desiderabile\fi & 
    #2 & 
    \ifnum#3=0 
        Piano di lavoro
    \else
        \ifnum#3=-1
            Intervista
        \else
            UC#3 §\ref{#3}
        \fi
    \fi \\
    \hline
}

% contatore RV
\newcounter{rvcount}
\setcounter{rvcount}{0}

\newcommand{\RV}[3]{%
    \stepcounter{rvcount}%
    RV-\thervcount & 
    \ifcase#1\or Obbligatorio\or Opzionale\or Desiderabile\fi & 
    #2 & 
    \ifnum#3=0 
        Piano di lavoro
    \else
        \ifnum#3=-1
            Intervista
        \else
            UC#3 §\ref{#3}
        \fi
    \fi \\
    \hline
}

\chapter{Analisi dei requisiti}
\label{cap:analisi-requisiti}

\intro{In questo capitolo vengono descritti in dettaglio tutti i requisiti funzionali, di qualità e di vincolo.}\\

\section{Casi d'uso}

Un caso d'uso è la descrizione dettagliata, tramite un diagramma UML e una descrizione testuale, di un insieme di scenari che hanno uno scopo comune per un attore all'interno del sistema.

% ---------------- ATTORI ----------------

\subsection{Attori principali}
L'attore principale è l'\textbf{utente}, suddiviso in utente base e admin.

\begingroup
\renewcommand{\arraystretch}{1.5}
\begin{center}
\begin{xltabular}{\textwidth}{|c|X|}
\hline
\textbf{Utente} & Utente generico registrato nel sistema \\
\hline
\textbf{Admin} & Amministratore di sistema \\
\hline
\textbf{Utente base} & Utente con permessi standard \\
\hline
\caption{Attori principali}
\end{xltabular}
\end{center}
\endgroup

\subsection{Attori secondari}
L'attore secondario è il modello di intelligenza artificiale (\textbf{LLM}), il quale viene utilizzato dal sistema per 
svolgere le funzionalità definite nei \gls{toolg} esposti, tra cui l'azione di \gls{crawlingg} dei vari siti web.
\begingroup
\renewcommand{\arraystretch}{1.5}
\begin{center}
\begin{xltabular}{\textwidth}{|c|X|}
\hline
\textbf{LLM} & Modello di intelligenza artificiale per l'elaborazione \\
\hline
\caption{Attori secondari}
\end{xltabular}
\end{center}
\endgroup

\section{Lista dei casi d'uso}

\subsection{UC1 - Autenticazione}
\label{1}
\makeUC{UC1}
{Utente}
{   \begin{itemize}
        \item Il sistema è attivo;
        \item l'utente si trova nella pagina di autenticazione.
    \end{itemize}
}
{L'utente è autenticato nel sistema.}
{   \begin{enumerate}
        \item L'utente si autentica con le proprie credenziali:
            \begin{itemize}
                \item inserimento \textit{username} (§\ref{1.1});
                \item inserimento \textit{password} (§\ref{1.2}).
            \end{itemize}
        \item l'utente conferma l'operazione.
    \end{enumerate}
}
{   \begin{itemize}
        \item UC1.1 §\ref{1.1}
        \item UC1.2 §\ref{1.2}
    \end{itemize} 
}
{L'utente non è registrato nel sistema o la \textit{password} è errata (§\ref{1.3}).}
{UC1.3 \ref{1.3}}{}


\subsection{UC1.1 - Inserimento \textit{username}}
\label{1.1}
\makeUC{}
{Utente}
{L'utente sta eseguendo l'autenticazione.}
{Il sistema riceve l'\textit{username} dell'utente.}
{L'utente inserisce l'\textit{username}.}{}{}{}{}


\subsection{UC1.2 - Inserimento \textit{password}}
\label{1.2}
\makeUC{}
{Utente}
{L'utente sta eseguendo l'autenticazione.}
{Il sistema riceve la \textit{password} dell'utente.}
{L'utente inserisce la \textit{password}.}{}{}{}{}

\subsection{UC1.3 - Visualizzazione errore \textit{username} o \textit{password} errati}
\label{1.3}
\makeUC{}
{Utente}
{L'utente ha inserito uno \textit{username} che non corrisponde ad un utente registrato nel sistema 
oppure la \textit{password} è errata per l'\textit{username} inserito.}
{Il sistema non autentica l'utente.}
{Il sistema mostra un messaggio di errore.}{}{}{}{}

\subsection{UC2 - \textit{Logout}}
\label{2}
\makeUC{UC2}
{Utente}
{
    \begin{itemize}
        \item Il sistema è attivo;
        \item l'utente è autenticato nel sistema;
        \item l'utente si trova nella pagina di gestione del profilo.
    \end{itemize}
}
{Il sistema effettua il \textit{logout} dell'utente, riportandolo alla pagina di autenticazione.}
{L'utente seleziona l'opzione relativa al \textit{logout}.}{}{}{}{}

\subsection{UC3 - Creazione utente}
\label{3}
\makeUC{UC3}
{Admin}
{   \begin{itemize}
        \item Il sistema è attivo;
        \item l'admin è autenticato nel sistema;
        \item l'admin si trova nella pagina di gestione degli utenti.
    \end{itemize}
}
{L'admin crea un nuovo utente base.}
{   \begin{enumerate}
        \item L'admin crea un nuovo utente base:
            \begin{itemize}
                \item inserimento \textit{email} (§\ref{3.1});
                \item inserimento \textit{username} (§\ref{3.2});
                \item inserimento \textit{password} (§\ref{3.3});
                \item selezione dei permessi utente (§\ref{3.4}).
            \end{itemize}
        \item l'admin conferma l'operazione.
    \end{enumerate}
}
{   \begin{itemize}
        \item UC3.1 §\ref{3.1}
        \item UC3.2 §\ref{3.2}
        \item UC3.3 §\ref{3.3}
        \item UC3.4 §\ref{3.4}
    \end{itemize} 
}
{
    \begin{itemize}
        \item L'\textit{email} è già esistente nel sistema (§\ref{3.5});
        \item l'\textit{username} è già esistente nel sistema (§\ref{3.6}).
    \end{itemize}
}
{
    \begin{itemize}
        \item UC3.5 §\ref{3.5}
        \item UC3.6 §\ref{3.6}
    \end{itemize}
}{}

\subsection{UC3.1 - Inserimento \textit{email}}
\label{3.1}
\makeUC{}
{Admin}
{L'admin sta eseguendo la creazione di un nuovo utente base.}
{Il sistema riceve l'\textit{email} del nuovo utente base.}
{L'admin inserisce l'\textit{email}.}{}{}{}{}

\subsection{UC3.2 - Inserimento \textit{username}}
\label{3.2}
\makeUC{}
{Admin}
{L'admin sta eseguendo la creazione di un nuovo utente base.}
{Il sistema riceve l'\textit{username} del nuovo utente base.}
{L'admin inserisce l'\textit{username}.}{}{}{}{}

\subsection{UC3.3 - Inserimento \textit{password}}
\label{3.3}
\makeUC{}
{Admin}
{L'admin sta eseguendo la creazione di un nuovo utente base.}
{Il sistema riceve la \textit{password} del nuovo utente base.}
{L'admin inserisce la \textit{password}.}{}{}{}{}

\subsection{UC3.4 - Selezione \textit{permessi}}
\label{3.4}
\makeUC{}
{Admin}
{L'admin sta eseguendo la creazione di un nuovo utente base.}
{Il sistema riceve la scelta dei permessi del nuovo utente base.}
{L'admin seleziona i permessi di utilizzo dei vari \gls{toolg} da assegnare all'utente. 
In particolare, per ogni tipo di azione disponibile, i permessi assegnabili sono:
\begin{itemize}
    \item utilizzo generale dell'azione: l'utente può chiamare e usare quel determinato \gls{toolg};
    \item utilizzo della cache: il sistema, quando avvia l'elaborazione per quel determinato \gls{toolg}, utilizza il controllo della \textit{cache};
    \item visualizzazione passaggi: il sistema, quando avvia l'elaborazione per quel determinato \gls{toolg}, mostra all'utente i passaggi di elaborazione.
\end{itemize}
}{}{}{}{}

\subsection{UC3.5 - Visualizzazione errore \textit{email} già esistente}
\label{3.5}
\makeUC{}
{Admin}
{L'admin ha inserito una \textit{email} che è già utilizzata all'interno del sistema.}
{Il sistema non registra l'utente base.}
{Il sistema mostra un messaggio di errore.}{}{}{}{}

\subsection{UC3.6 - Visualizzazione errore \textit{username} già esistente}
\label{3.6}
\makeUC{}
{Admin}
{L'admin ha inserito uno \textit{username} che è già utilizzato all'interno del sistema.}
{Il sistema non registra l'utente base.}
{Il sistema mostra un messaggio di errore.}{}{}{}{}

\subsection{UC4 - Eliminazione utente base}
\label{4}
\makeUC{}
{Admin}
{   \begin{itemize}
        \item Il sistema è attivo;
        \item l'admin è autenticato nel sistema;
        \item l'admin si trova nella pagina di gestione degli utenti.
    \end{itemize}
}
{Il sistema elimina il profilo dell'utente base selezionato.}
{   \begin{enumerate}
        \item l'admin seleziona l'utente base che vuole eliminare;
        \item l'admin seleziona l'opzione di eliminazione dell'utente base;
        \item l'admin visualizza un pop-up di conferma di eliminazione del profilo, quindi conferma.
    \end{enumerate}
}
{}{}{}{}

\subsection{UC5 - Visualizzazione profilo}
\label{5}
\makeUC{UC5}
{Utente}
{   \begin{itemize}
        \item Il sistema è attivo;
        \item l'utente è autenticato nel sistema;
        \item l'utente si trova nella pagina di gestione del profilo.
    \end{itemize}
}
{L'utente visualizza il proprio profilo.}
{L'utente visualizza le informazioni relative al profilo:
    \begin{itemize}
        \item visualizzazione \textit{username} (§\ref{5.1});
        \item visualizzazione \textit{email} (§\ref{5.2}).
    \end{itemize}
}
{   \begin{itemize}
        \item UC5.1 §\ref{5.1}
        \item UC5.2 §\ref{5.2}
    \end{itemize} 
}{}{}{}

\subsection{UC5.1 - Visualizzazione \textit{username}}
\label{5.1}
\makeUC{}
{Utente}
{L'utente sta visualizzando il proprio profilo.}
{L'utente visualizza l'\textit{username}.}
{L'utente visualizza l'\textit{username}.}{}{}{}{}

\subsection{UC5.2 - Visualizzazione \textit{email}}
\label{5.2}
\makeUC{}
{Utente}
{L'utente sta visualizzando il proprio profilo.}
{L'utente visualizza l'\textit{email}.}
{L'utente visualizza l'\textit{email}.}{}{}{}{}

\subsection{UC6 - Modifica profilo}
\label{6}
\makeUC{UC6}
{Utente}
{   \begin{itemize}
        \item Il sistema è attivo;
        \item l'utente è autenticato nel sistema;
        \item l'utente si trova nella pagina di gestione del profilo.
    \end{itemize}
}
{L'utente modifica le proprie credenziali.}
{   \begin{enumerate}
        \item L'utente seleziona l'opzione di modifica del profilo;
        \item l'utente modifica le credenziali:
            \begin{itemize}
                \item modifica \textit{username} (§\ref{6.1});
                \item modifica \textit{password} (§\ref{6.2}).
            \end{itemize}
        \item l'utente conferma l'operazione.
    \end{enumerate}
}
{   \begin{itemize}
        \item UC6.1 §\ref{6.1}
        \item UC6.2 §\ref{6.2}
    \end{itemize} 
}
{
    \begin{itemize}
        \item l'\textit{username} è già esistente nel sistema (§\ref{3.6}).
    \end{itemize}
}
{
    \begin{itemize}
        \item UC3.6 §\ref{3.6}
    \end{itemize}
}{}

\subsection{UC6.1 - Modifica \textit{username}}
\label{6.1}
\makeUC{}
{Utente}
{L'utente sta eseguendo la modifica del profilo.}
{Il sistema riceve l'\textit{username} dell'utente.}
{L'utente inserisce l'\textit{username}.}{}{}{}{}

\subsection{UC6.2 - Modifica \textit{password}}
\label{6.2}
\makeUC{}
{Utente}
{L'utente sta eseguendo la modifica del profilo.}
{Il sistema riceve la \textit{password} dell'utente.}
{L'utente inserisce la \textit{password}.}{}{}{}{}

\subsection{UC7 - Invio messaggio in linguaggio naturale}
\label{7}
\makeUC{UC7}
{Utente, LLM}
{   \begin{itemize}
        \item Il sistema è attivo;
        \item l'utente è autenticato nel sistema;
        \item l'utente si trova nella pagina relativa alla \textit{chat}.
    \end{itemize}
}
{Il sistema riceve il messaggio dell'utente, comprende il \gls{toolg} richiesto, 
estrae i parametri necessari alla sua esecuzione e l'\gls{llmg} avvia l'elaborazione.}
{   \begin{enumerate}
        \item L'utente definisce, in linguaggio naturale, l'azione da intraprendere
        (per esempio ''\textit{cerca qual è il meteo attuale a Padova nel sito meteo.it. Rispondi in formato Markdown.}'');
        \item l'utente definisce, in linguaggio naturale, i parametri di ricerca:
            \begin{itemize}
                \item inserimento url (§\ref{7.3});
                \item inserimento domanda (§\ref{7.4});
                \item inserimento tipo di formato della risposta (§\ref{7.5}) (opzionale).
            \end{itemize}
        \item in base ai permessi assegnati all'utente, la ricerca può essere effettuata:
            \begin{itemize}
                \item con controllo della \textit{cache} per evitare richieste ridondanti (§\ref{7.1});
                \item senza controllo della \textit{cache} (\textit{cache bypass}) (§\ref{7.2}).
            \end{itemize}
        \item l'utente invia il messaggio.
    \end{enumerate}
}
{   \begin{itemize}
        \item UC7.3 §\ref{7.3}
        \item UC7.4 §\ref{7.4}
        \item UC7.5 §\ref{7.5}
    \end{itemize} 
}
{
    \begin{itemize}
        \item Alcuni parametri obbligatori per l'esecuzione del \gls{toolg} sono mancanti (§\ref{7.6});
        \item L'azione richiesta dall'utente non è permessa (§\ref{7.7}).
    \end{itemize}
}
{
    \begin{itemize}
        \item UC7.6 §\ref{7.6}
        \item UC7.7 §\ref{7.7}
    \end{itemize}
}{
    \begin{itemize}
        \item UC7.1 §\ref{7.1}
        \item UC7.2 §\ref{7.2}
    \end{itemize}
}

\subsection{UC7.1 - Avvio tool con controllo \textit{cache}}
\label{7.1}
\makeUC{}
{Utente, LLM}
{
    \begin{itemize}
        \item L'utente si trova nella pagina relativa alla \textit{chat};
        \item L'utente ha inviato un messaggio richiedendo un \gls{toolg};
        \item L'utente ha i permessi di utilizzo del \gls{toolg} e i permessi di uso della \textit{cache}.
    \end{itemize}
}
{Il sistema avvia l'azione in base al messaggio inviato dall'utente.}
{L'\gls{llmg} avvia l'elaborazione.}{}{}{}{}

\subsection{UC7.2 - Avvio tool con \textit{cache bypass}}
\label{7.2}
\makeUC{}
{Utente, LLM}
{
    \begin{itemize}
        \item L'utente si trova nella pagina relativa alla \textit{chat};
        \item L'utente ha inviato un messaggio richiedendo un \gls{toolg};
        \item L'utente ha i permessi di utilizzo del \gls{toolg} ma non ha i permessi di uso della \textit{cache}.
    \end{itemize}
}
{Il sistema avvia correttamente l'azione con \textit{cache bypass}, ovvero non effettua il 
controllo della \textit{cache}.}
{L'\gls{llmg} avvia l'elaborazione.}{}{}{}{}

\subsection{UC7.3 - Inserimento url}
\label{7.3}
\makeUC{}
{Utente}
{L'utente sta inviando un messaggio nella pagina di \textit{chat}.}
{Il sistema riceve l'url.}
{L'utente inserisce l'url.}{}{}{}{}

\subsection{UC7.4 - Inserimento domanda}
\label{7.4}
\makeUC{}
{Utente}
{L'utente sta inviando un messaggio nella pagina di \textit{chat}.}
{Il sistema riceve la domanda.}
{L'utente inserisce la domanda.}{}{}{}{}

\subsection{UC7.5 - Inserimento tipo di formato della risposta}
\label{7.5}
\makeUC{}
{Utente}
{L'utente sta inviando un messaggio nella pagina di \textit{chat}.}
{Il sistema riceve il tipo di formato con cui generare la risposta.}
{L'utente inserisce il tipo di formato della risposta.}{}{}{}{}

\subsection{UC7.6 - Visualizzazione errore parametri mancanti}
\label{7.6}
\makeUC{}
{Utente}
{L'utente non ha inserito uno o più parametri obbligatori per l'elaborazione del \gls{toolg} richiesto.}
{Il sistema non avvia il \gls{toolg}}
{Il sistema mostra un messaggio di errore e 
segnala all'utente quali sono i parametri mancanti.}{}{}{}{}

\subsection{UC7.7 - Visualizzazione errore tool non permesso}
\label{7.7}
\makeUC{}
{Utente}
{L'utente non ha i permessi necessari per eseguire l'azione richiesta 
oppure il \gls{toolg} richiesto non è valido.}
{Il sistema non avvia l'elaborazione.}
{Il sistema mostra un messaggio di errore.}{}{}{}{}

\subsection{UC8 - Visualizzazione lista dei passaggi della ricerca}
\label{8}
\makeUC{}
{Utente, LLM}
{   \begin{itemize}
        \item Il sistema è attivo;
        \item l'utente ha inviato un messaggio nella pagina relativa alla \textit{chat};
        \item l'LLM sta elaborando o ha elaborato la richiesta dell'utente;
        \item l'utente ha i permessi per visualizzare i passaggi della ricerca.
    \end{itemize}
}
{L'utente visualizza la sequenza di azioni logiche compiute dall'\gls{llmg}.}
{L'utente visualizza in tempo reale i singoli passaggi generati dall'attività dell'\gls{llmg} (§\ref{8.1}).}
{
    \begin{itemize}
        \item UC8.1 §\ref{8.1}
    \end{itemize}
}
{}{}{}

\subsection{UC8.1 - Visualizzazione singolo passaggio della ricerca}
\label{8.1}
\makeUC{}
{Utente, LLM}
{   \begin{itemize}
        \item L'utente ha inviato un messaggio nella pagina relativa alla \textit{chat};
        \item l'utente ha i permessi per visualizzare i passaggi della ricerca;
        \item l'utente sta visualizzando il flusso di ricerca generato dall'\gls{llmg}.
    \end{itemize}
}
{L'utente visualizza un messaggio informativo riguardante l'operazione specifica effettuata dall'\gls{llmg} in quel momento.}
{
    \begin{enumerate}
        \item L'utente visualizza una delle seguenti tipologie di passaggi, in caso il \gls{toolg} sia quello di \textit{crawling}:
            \begin{itemize}
                \item pagina web estratta e analizzata (§\ref{8.1.1}); 
                \item riassunto semantico del contenuto (§\ref{8.1.2}); 
                \item link trovati nella pagina web analizzata (§\ref{8.1.3}); 
                \item link selezionati per continuare il processo di \textit{crawling} (§\ref{8.1.4}); 
                \item risposta finale elaborata (§\ref{8.1.5}). 
            \end{itemize}
        \item L'utente, per il singolo passaggio, visualizza le seguenti proprietà:
            \begin{itemize}
                \item il nome del passaggio (§\ref{8.2});
                \item la descrizione del passaggio (§\ref{8.3}).
            \end{itemize}
        \item Se al termine dell'elaborazione il sistema genera correttamente la risposta,
        l'utente può salvare la risposta generata (§\ref{8.5}).
    \end{enumerate}
}
{
    \begin{itemize}
        \item UC8.2 §\ref{8.2}
        \item UC8.3 §\ref{8.3}
    \end{itemize}
}
{Interruzione dell'elaborazione per errore interno o di connessione (§\ref{8.4}).}
{
    \begin{itemize}
        \item UC8.4 §\ref{8.4}
    \end{itemize}
}
{
    \begin{itemize}
        \item UC8.1.1 §\ref{8.1.1}
        \item UC8.1.2 §\ref{8.1.2}
        \item UC8.1.3 §\ref{8.1.3}
        \item UC8.1.4 §\ref{8.1.4}
        \item UC8.1.5 §\ref{8.1.5}
    \end{itemize}
}

\subsection{UC8.1.1 - Visualizzazione passaggio: pagina analizzata}
\label{8.1.1}
\makeUC{}
{Utente, LLM}
{L'LLM sta estraendo il contenuto da una pagina \textit{web}.}
{L'utente visualizza il nome della pagina \textit{web} da cui si sta estraendo il contenuto.}
{Il sistema mostra un messaggio informativo relativo all'estrazione del contenuto da una pagina \textit{web}.}{}{}{}{}

\subsection{UC8.1.2 - Visualizzazione passaggio: riassunto della pagina analizzata}
\label{8.1.2}
\makeUC{}
{Utente, LLM}
{L'LLM ha completato l'estrazione del testo da una pagina \textit{web}.}
{L'utente visualizza la sintesi generata dall'LLM.}
{Il sistema mostra il riassunto informativo che l'LLM ha estratto dal link analizzato per determinare la pertinenza alla domanda.}{}{}{}{}

\subsection{UC8.1.3 - Visualizzazione passaggio: link trovati}
\label{8.1.3}
\makeUC{}
{Utente}
{
    \begin{itemize}
        \item Il sistema ha estratto il testo riassuntivo da una pagina \textit{web};
        \item il sistema non ha ritenuto sufficiente o pertinente il contenuto trovato per poter rispondere in maniera completa alla domanda dell'utente;
        \item il sistema ha completato l'estrazione di tutti i link a cui fa riferimento la pagina \textit{web} per poter continuare la ricerca.
    \end{itemize}
}
{L'utente visualizza la lista dei link trovati.}
{Il sistema mostra la lista dei link trovati nella pagina \textit{web} analizzata.}{}{}{}{}

\subsection{UC8.1.4 - Visualizzazione passaggio: link scelti}
\label{8.1.4}
\makeUC{}
{Utente, LLM}
{L'LLM ha completato la scelta dei link pertinenti per la ricerca.}
{L'utente visualizza la lista dei link scelti dall'LLM per continuare la ricerca.
Il massimo numero di link scelti corrisponde al parametro "massima ampiezza" 
definito dal \gls{llmg}.}
{Il sistema mostra la lista dei link scelti.}{}{}{}{}

\subsection{UC8.1.5 - Visualizzazione passaggio: risposta}
\label{8.1.5}
\makeUC{}
{Utente, LLM}
{
    \begin{itemize}
        \item Il sistema ha estratto il testo riassuntivo da una pagina \textit{web};
        \item il sistema ha ritenuto sufficiente e pertinente il contenuto trovato per poter rispondere all'utente.
    \end{itemize}
}
{L'utente visualizza il passaggio di risposta.}
{
    \begin{enumerate}
        \item L'LLM genera la risposta finale e la mostra all'utente;
        \item il sistema termina la ricerca.
    \end{enumerate}
}{}{}{}{}

\subsection{UC8.2 - Visualizzazione nome passaggio}
\label{8.2}
\makeUC{}
{Utente}
{L'utente sta visualizzando i passaggi della ricerca avviata.}
{L'utente visualizza il nome del passaggio.}
{Il sistema mostra il nome relativo al passaggio visualizzato.}{}{}{}{}

\subsection{UC8.3 - Visualizzazione descrizione passaggio}
\label{8.3}
\makeUC{}
{Utente}
{L'utente sta visualizzando i passaggi della ricerca avviata.}
{L'utente visualizza la descrizione del passaggio.}
{Il sistema mostra la descrizione relativa al passaggio visualizzato.}{}{}{}{}

\subsection{UC8.4 - Visualizzazione errore durante la ricerca}
\label{8.4}
\makeUC{}
{Utente, LLM}
{L'LLM riscontra un'impossibilità nel proseguire l'elaborazione della risposta.}
{Il sistema notifica il fallimento della ricerca e ne interrompe l'esecuzione.}
{Il sistema mostra un messaggio di errore indicando se l'interruzione è dovuta a problemi interni (es. \textit{timeout}, limiti di contesto) o problemi di rete.}{}{}{}{}

\subsection{UC8.5 - Salva risposta}
\label{8.5}
\makeUC{UC8.5}
{Utente}
{L'utente sta visualizzando la risposta generata correttamente dal sistema.}
{Il contenuto della risposta viene correttamente salvato.}
{
    L'utente seleziona una tra le seguenti opzioni:
        \begin{itemize}
            \item Copia la risposta negli appunti (§\ref{8.5.1});
            \item Salva in un file la risposta (§\ref{8.5.2}).
        \end{itemize}
}{}{}{}{
    \begin{itemize}
        \item UC8.5.1 §\ref{8.5.1}
        \item UC8.5.2 §\ref{8.5.2}
    \end{itemize}
}

\subsection{UC8.5.1 - Copia risposta negli appunti}
\label{8.5.1}
\makeUC{}
{Utente}
{L'utente sta visualizzando la risposta generata correttamente dal sistema.}
{Il contenuto della risposta viene correttamente copiato negli appunti.}
{L'utente seleziona l'opzione di copia negli appunti.}{}{}{}{}

\subsection{UC8.5.2 - Salvataggio risposta in un file}
\label{8.5.2}
\makeUC{}
{Utente}
{L'utente sta visualizzando la risposta generata correttamente dal sistema.}
{Il contenuto della risposta viene correttamente salvato in un file nel dispositivo dell'utente.}
{
    \begin{enumerate}
        \item L'utente seleziona l'opzione di creazione di un file in cui salvare la risposta;
        \item l'utente sceglie il nome e la posizione in cui salvare il file;
        \item l'utente conferma l'operazione.
    \end{enumerate}
}{}{
    Il file non viene creato e salvato a causa di un errore durante il procedimento (§\ref{8.5.2.1}).
}{
    \begin{itemize}
        \item UC8.5.2.1 §\ref{8.5.2.1}
    \end{itemize}
}{}

\subsection{UC8.5.2.1 - Visualizzazione errore file non salvato}
\label{8.5.2.1}
\makeUC{}
{Utente}
{
    \begin{itemize}
        \item L'utente sta visualizzando la risposta generata correttamente dal sistema;
        \item l'utente ha selezionato l'opzione di salvataggio della risposta in un file.
    \end{itemize}
}
{Il contenuto della risposta non viene salvato in un file.}
{Il sistema mostra all'utente un messaggio di errore.}{}{}{}{}

\subsection{UC9 - Visualizzazione storico delle chat}
\label{9}
\makeUC{}
{Utente}
{
    \begin{itemize}
        \item Il sistema è attivo;
        \item l'utente è autenticato nel sistema;
        \item l'utente sta visualizzando lo storico delle \textit{chat}.
    \end{itemize}
}
{
    L'utente visualizza la lista di tutte le ricerche nel periodo selezionato.
}
{
    \begin{enumerate}
        \item L'utente seleziona un filtro di periodo, per esempio ''\textit{ultima settimana}, \textit{ultimi 30 giorni}, \dots'' (§\ref{9.1});
        \item l'utente conferma l'operazione;
        \item l'utente visualizza i singoli elementi dello storico (§\ref{9.2}).
    \end{enumerate}
}
{
    \begin{itemize}
        \item UC9.1 \ref{9.1}
        \item UC9.2 \ref{9.2}
    \end{itemize}
}
{}
{}
{}

\subsection{UC9.1 - Selezione filtro periodo}
\label{9.1}
\makeUC{}
{Utente}
{L'utente sta visualizzando lo storico delle \textit{chat}.}
{L'utente visualizza lo storico delle \textit{chat} nel periodo selezionato.}
{Il sistema mostra la lista di tutte le \textit{chat} nel periodo selezionato 
dall'utente.}{}{}{}{}

\subsection{UC9.2 - Visualizzazione singolo elemento dello storico}
\label{9.2}
\makeUC{}
{Utente}
{L'utente sta visualizzando lo storico delle \textit{chat}.}
{L'utente visualizza il singolo elemento dello storico.}
{Il sistema mostra il singolo elemento dello storico, caratterizzato dall'ultimo messaggio 
inviato nella \textit{chat}.
}{
    \begin{itemize}
        \item UC9.2.1 §\ref{9.2.1}
    \end{itemize}
}{}{}{}

\subsection{UC9.2.1 - Visualizzazione ultimo messaggio della chat}
\label{9.2.1}
\makeUC{}
{Utente}
{L'utente sta visualizzando lo storico delle \textit{chat} nel periodo selezionato.}
{L'utente visualizza l'ultimo messaggio della \textit{chat} che può essere un messaggio dell'utente 
oppure un messaggio del sistema.}
{Il sistema mostra parte dell'ultimo messaggio della \textit{chat} nello storico.}{}{}{}{}


\subsection{UC10 - Eliminazione storico}
\label{10}
\makeUC{UC10}
{Utente}
{
    \begin{itemize}
        \item Il sistema è attivo;
        \item l'utente è autenticato nel sistema;
        \item l'utente si trova nella pagina di storico delle ricerche;
        \item l'utente ha i permessi per poter visualizzare lo storico delle ricerche.
    \end{itemize}
}
{
    L'utente elimina correttamente le richieste dallo storico.
}
{
    L'utente può effettuare una delle seguenti tipologie di eliminazione dallo storico:
    \begin{itemize}
        \item Eliminazione del singolo elemento dello storico (§\ref{10.1});
        \item eliminazione di tutto lo storico (§\ref{10.2}).
    \end{itemize}
}
{}
{
    Lo storico è vuoto (§\ref{10.3}).
}
{
    \begin{itemize}
        \item UC10.3 §\ref{10.3}
    \end{itemize}
}
{
    \begin{itemize}
        \item UC10.1 §\ref{10.1}
        \item UC10.2 §\ref{10.2}
    \end{itemize}
}

\subsection{UC10.1 - Eliminazione singolo elemento dello storico}
\label{10.1}
\makeUC{}
{Utente}
{L'utente sta visualizzando lo storico delle \textit{chat} nel periodo selezionato.}
{L'utente elimina correttamente il singolo elemento dello storico selezionato.}
{
    \begin{enumerate}
        \item L'utente seleziona l'elemento che vuole eliminare;
        \item l'utente seleziona l'opzione relativa all'eliminazione dell'elemento;
        \item l'utente conferma l'operazione.
    \end{enumerate}
}{}{}{}{}

\subsection{UC10.2 - Eliminazione intero storico}
\label{10.2}
\makeUC{}
{Utente}
{L'utente sta visualizzando lo storico delle \textit{chat} nel periodo selezionato.}
{L'utente elimina correttamente tutto lo storico, indipendentemente dal periodo selezionato.}
{
    \begin{enumerate}
        \item l'utente seleziona l'opzione relativa all'eliminazione di tutto lo storico, indipendentemente dal periodo selezionato;
        \item l'utente conferma l'operazione.
    \end{enumerate}
}{}{}{}{}

\subsection{UC10.3 - Visualizzazione errore storico vuoto}
\label{10.3}
\makeUC{}
{Utente}
{L'utente ha selezionato l'opzione di eliminazione dello storico delle richieste che è vuoto.}
{Il sistema non elimina lo storico.}
{Il sistema mostra un messaggio di errore.}{}{}{}{}

\subsection{Funzionali}
\begingroup
\renewcommand{\arraystretch}{1.50}
\begin{center}
    \begin{xltabular}{\textwidth}{|c|c| >{\arraybackslash}X|c|}
        \hline
        \textbf{ID} & \textbf{Tipo} & \textbf{Descrizione} & \textbf{Fonti} \\
        \hline
        % ── OBBLIGATORI ──────────────────────────────────────────────────────
        \makeReq{1}{L'utente deve potersi autenticare presso il sistema}{1}
        \makeReq{1}{L'utente deve inserire il proprio \textit{username} per autenticarsi presso il sistema}{1.1}
        \makeReq{1}{L'utente deve inserire la propria \textit{password} per autenticarsi presso il sistema}{1.2}
        \makeReq{1}{L'utente deve ricevere un errore se le credenziali inserite sono errate}{1.3}
        \makeReq{1}{L'utente deve poter effettuare il \textit{logout} dal sistema}{2}
        \makeReq{1}{L'admin deve poter creare un nuovo utente base nel sistema}{3}
        \makeReq{1}{L'admin deve poter inserire una \textit{email} per il nuovo utente base}{3.1}
        \makeReq{1}{L'admin deve poter inserire uno \textit{username} per il nuovo utente base}{3.2}
        \makeReq{1}{L'admin deve poter inserire una \textit{password} per il nuovo utente base}{3.3}
        \makeReq{1}{L'admin deve poter selezionare i permessi per il nuovo utente base}{3.4}
        \makeReq{1}{L'admin deve ricevere un errore se la \textit{email} inserita è già esistente nel sistema}{3.5}
        \makeReq{1}{L'admin deve ricevere un errore se lo \textit{username} inserito è già esistente nel sistema}{3.6}
        \makeReq{1}{L'admin deve poter eliminare un utente base dal sistema}{4}
        \makeReq{1}{L'utente deve poter visualizzare il proprio profilo}{5}
        \makeReq{1}{L'utente deve poter visualizzare il proprio \textit{username} nel profilo}{5.1}
        \makeReq{1}{L'utente deve poter visualizzare la propria \textit{email} nel profilo}{5.2}
        \makeReq{1}{L'utente deve poter modificare il proprio profilo}{6}
        \makeReq{1}{L'utente deve poter modificare il proprio \textit{username}}{6.1}
        \makeReq{1}{L'utente deve poter modificare la propria \textit{password}}{6.2}
        \makeReq{1}{L'utente deve poter inviare un messaggio in linguaggio naturale per avviare una ricerca}{7}
        \makeReq{1}{L'utente deve poter avviare una ricerca con \textit{cache bypass}}{7.2}
        \makeReq{1}{Il sistema deve estrarre l'URL di partenza dal messaggio dell'utente per avviare la ricerca}{7.3}
        \makeReq{1}{Il sistema deve estrarre la domanda dal messaggio dell'utente per avviare la ricerca}{7.4}
        \makeReq{1}{Il sistema deve segnalare all'utente i parametri obbligatori mancanti, impedendo l'avvio del \textit{tool}}{7.6}
        \makeReq{1}{L'utente deve ricevere un errore se l'azione richiesta non è permessa o il \textit{tool} richiesto non è valido}{7.7}
        \makeReq{1}{L'utente deve poter visualizzare la risposta generata dal sistema}{8.1.5}
        \makeReq{1}{Il sistema deve notificare l'utente in caso di errore durante la ricerca}{8.4}
        % ── OPZIONALI ─────────────────────────────────────────────────────────
        \makeReq{2}{L'utente deve poter specificare il formato della risposta (\textit{Markdown}, \textit{HTML}, \textit{JSON})}{7.5}
        \makeReq{2}{L'utente deve poter visualizzare in tempo reale i passaggi della ricerca in corso}{8}
        \makeReq{2}{L'utente deve poter visualizzare in tempo reale il singolo passaggio della ricerca in corso}{8.1}
        \makeReq{2}{L'utente deve poter visualizzare il passaggio relativo alla pagina \textit{web} in fase di analisi}{8.1.1}
        \makeReq{2}{L'utente deve poter visualizzare il passaggio relativo al riassunto semantico della pagina analizzata}{8.1.2}
        \makeReq{2}{L'utente deve poter visualizzare il passaggio relativo ai link trovati nella pagina analizzata}{8.1.3}
        \makeReq{2}{L'utente deve poter visualizzare il passaggio relativo ai link scelti per continuare la ricerca}{8.1.4}
        \makeReq{2}{L'utente deve poter visualizzare il nome del passaggio}{8.2}
        \makeReq{2}{L'utente deve poter visualizzare la descrizione del passaggio}{8.3}
        \makeReq{2}{L'utente deve poter visualizzare lo storico delle proprie \textit{chat}}{9}
        \makeReq{2}{L'utente deve poter filtrare lo storico per periodo temporale (ultima settimana, ultimo mese, \dots)}{9.1}
        \makeReq{2}{L'utente deve poter visualizzare i singoli elementi dello storico}{9.2}
        \makeReq{2}{L'utente deve poter visualizzare, per ogni elemento dello storico, l'ultimo messaggio della relativa \textit{chat}}{9.2.1}
        \makeReq{2}{L'utente deve poter eliminare lo storico delle proprie \textit{chat}}{10}
        \makeReq{2}{L'utente deve poter eliminare il singolo elemento dello storico}{10.1}
        \makeReq{2}{L'utente deve poter eliminare tutto lo storico}{10.2}
        \makeReq{2}{L'utente deve ricevere un errore se tenta di eliminare lo storico quando è vuoto}{10.3}
        % ── DESIDERABILI ──────────────────────────────────────────────────────
        \makeReq{3}{L'utente deve poter avviare una ricerca con controllo della \textit{cache}}{7.1}
        \makeReq{3}{L'utente deve poter salvare la risposta generata dal sistema}{8.5}
        \makeReq{3}{L'utente deve poter copiare la risposta generata negli appunti}{8.5.1}
        \makeReq{3}{L'utente deve poter salvare la risposta generata in un file locale}{8.5.2}
        \makeReq{3}{L'utente deve ricevere un errore se il file con la risposta generata non viene salvato correttamente}{8.5.2.1}
    \end{xltabular}
\end{center}
\endgroup

\subsection{Riepilogo}
\begingroup
\renewcommand{\arraystretch}{1.50}
\begin{center}
    \begin{xltabular}{\textwidth}{| c | c | c | c | c |}
        \hline
        \textbf{Tipi di requisiti} & \textbf{Obbligatori} & \textbf{Opzionali} & \textbf{Desiderabili} & \textbf{Totali} \\
        \hline
        \textbf{Funzionali} & 27 & 17 & 5 & 49 \\
        \hline
        \textbf{Di qualità} & 6 & 0 & 0 & 6 \\
        \hline
        \textbf{Di vincolo} & 2 & 0 & 0 & 2 \\
        \hline
        \textbf{Totali} & 35 & 17 & 5 & \textbf{57} \\
        \hline
    \end{xltabular}
\end{center}
\endgroup